\documentclass[12pt]{article}
\usepackage[utf8]{inputenc}
\usepackage[T1]{fontenc}
\usepackage{hyperref}
\title{Ayautla Mazatec Language Resources for K-12 Curriculum Development}
\author{Compiled by Claude Sonnet 5 (claude-sonnet-5) for Andr\'as Kornai}
\date{}
\begin{document}
\maketitle

\section{Introduction}
Ayautla Mazatec (ISO 639-3 vmy; Glottocode ayau1235) is an Otomanguean language of Oaxaca, Mexico, with about 3,700 speakers. Reported figures include roughly 2,800 monolinguals, and it is reported to be about 80\% intelligible with Huautla Mazatec. The search found little beyond Scripture and audio materials. No school materials, grammar or dedicated dictionary turned up.

\section{Primary Educational Resources}
No primers, textbooks or school programmes for Ayautla Mazatec were found in the search results.

\section{Linguistic and Reference Materials}
\begin{itemize}
\item Glottolog entry for Ayautla Mazatec, with classification and alternative names; its bibliography may list further work (not checked): \url{https://glottolog.org/resource/languoid/id/ayau1235}
\item Wikipedia article on Ayautla Mazatec, with speaker figures: \url{https://en.wikipedia.org/wiki/Ayautla_Mazatec}
\item Glosbe online Ayautla Mazatec--German dictionary (content and depth not verified): \url{https://glosbe.com/vmy/de}
\end{itemize}

\section{Religious and Community Resources}
\begin{itemize}
\item New Testament in Ayautla Mazatec (2009, Wycliffe), the first Scripture translation: \url{https://dbs.org/bibles/VMYWBT}
\item Scripture Earth page with readable New Testament, audio, video and apps: \url{https://www.scriptureearth.org/00eng.php?iso=vmy}
\end{itemize}

\section{Digital and Audio Resources}
\begin{itemize}
\item Global Recordings Network audio recordings (evangelism and basic Bible teaching): \url{https://globalrecordings.net/en/language/vmy}
\end{itemize}

\section{Curriculum Development Context}
The language is stable and used as a first language across the community, but no school use is known. Any curriculum would need new material. The New Testament and audio could serve as a starting corpus for orthography and reading texts, and a community-led effort would be needed for the rest. The closely related Huautla Mazatec may offer more materials.
\end{document}